\section{Quasi-coherence of higher direct images}
\label{section-quasi-coherence}

\noindent
We have seen that the higher cohomology groups of a quasi-coherent module on
an affine are zero. For (quasi-)separated quasi-compact schemes $X$ this implies
vanishing of cohomology groups of quasi-coherent sheaves beyond a certain
degree. However, it may not be the case that $X$ has finite cohomological
dimension, because that is defined in terms of vanishing of cohomology
of {\it all} $\mathcal{O}_X$-modules.

\begin{lemma}[Induction Principle]
\label{lemma-induction-principle}
\begin{reference}
\cite[Proposition 3.3.1]{BvdB}
\end{reference}
Let $X$ be a quasi-compact and quasi-separated scheme. Let $P$ be a property
of the quasi-compact opens of $X$. Assume that
\begin{enumerate}
\item $P$ holds for every affine open of $X$,
\item if $U$ is quasi-compact open, $V$ affine open,
$P$ holds for $U$, $V$, and $U \cap V$, then
$P$ holds for $U \cup V$.
\end{enumerate}
Then $P$ holds for every quasi-compact open of $X$
and in particular for $X$.
\end{lemma}

\begin{proof}
First we argue by induction that $P$ holds for {\it separated} quasi-compact
opens $W \subset X$. Namely, such an open can be written as
$W = U_1 \cup \ldots \cup U_n$ and we can do induction on $n$ using
property (2) with $U = U_1 \cup \ldots \cup U_{n - 1}$ and $V = U_n$.
This is allowed because
$U \cap V = (U_1 \cap U_n) \cup \ldots \cup (U_{n - 1} \cap U_n)$
is also a union of $n - 1$ affine open subschemes by
Schemes, Lemma \ref{schemes-lemma-characterize-separated}
applied to the affine opens $U_i$ and $U_n$ of $W$.
Having said this, for any quasi-compact open $W \subset X$ we can
do induction on the number of affine opens needed to cover $W$
using the same trick as before and using that the quasi-compact open
$U_i \cap U_n$ is separated as an open subscheme of the affine scheme $U_n$.
\end{proof}

\begin{lemma}
\label{lemma-vanishing-nr-affines}
\begin{slogan}
For schemes with affine diagonal, the cohomology of quasi-coherent
modules vanishes in degrees bigger than the number of affine
opens needed in a covering.
\end{slogan}
Let $X$ be a quasi-compact scheme with affine diagonal (for example
if $X$ is separated).
Let $t = t(X)$ be the minimal number of affine opens needed to
cover $X$. Then $H^n(X, \mathcal{F}) = 0$ for all $n \geq t$ and all
quasi-coherent sheaves $\mathcal{F}$.
\end{lemma}

\begin{proof}
First proof.
By induction on $t$.
If $t = 1$ the result follows from
Lemma \ref{lemma-quasi-coherent-affine-cohomology-zero}.
If $t > 1$ write $X = U \cup V$ with $V$ affine open and
$U = U_1 \cup \ldots \cup U_{t - 1}$ a union of $t - 1$ open affines.
Note that in this case
$U \cap V =  (U_1 \cap V) \cup \ldots (U_{t - 1} \cap V)$
is also a union of $t - 1$ affine open subschemes.
Namely, since the diagonal is affine, the intersection of two
affine opens is affine, see Lemma \ref{lemma-affine-diagonal}.
We apply the Mayer-Vietoris long exact sequence
$$
0 \to
H^0(X, \mathcal{F}) \to
H^0(U, \mathcal{F}) \oplus H^0(V, \mathcal{F}) \to
H^0(U \cap V, \mathcal{F}) \to
H^1(X, \mathcal{F}) \to \ldots
$$
see Cohomology, Lemma \ref{cohomology-lemma-mayer-vietoris}.
By induction we see that the groups $H^i(U, \mathcal{F})$,
$H^i(V, \mathcal{F})$, $H^i(U \cap V, \mathcal{F})$ are zero for
$i \geq t - 1$. It follows immediately that $H^i(X, \mathcal{F})$
is zero for $i \geq t$.

\medskip\noindent
Second proof.
Let $\mathcal{U} : X = \bigcup_{i = 1}^t U_i$ be a finite affine open
covering. Since $X$ is has affine diagonal the multiple intersections
$U_{i_0 \ldots i_p}$ are all affine, see
Lemma \ref{lemma-affine-diagonal}.
By Lemma \ref{lemma-cech-cohomology-quasi-coherent} the {\v C}ech
cohomology groups $\check{H}^p(\mathcal{U}, \mathcal{F})$
agree with the cohomology groups. By
Cohomology, Lemma \ref{cohomology-lemma-alternating-usual}
the {\v C}ech cohomology groups may be computed using the alternating
{\v C}ech complex $\check{\mathcal{C}}_{alt}^\bullet(\mathcal{U}, \mathcal{F})$.
As the covering consists of $t$ elements we see immediately
that $\check{\mathcal{C}}_{alt}^p(\mathcal{U}, \mathcal{F}) = 0$
for all $p \geq t$. Hence the result follows.
\end{proof}

\begin{lemma}
\label{lemma-affine-diagonal-universal-delta-functor}
Let $X$ be a quasi-compact scheme with affine diagonal
(for example if $X$ is separated). Then
\begin{enumerate}
\item given a quasi-coherent $\mathcal{O}_X$-module $\mathcal{F}$
there exists an embedding $\mathcal{F} \to \mathcal{F}'$ of
quasi-coherent $\mathcal{O}_X$-modules
such that $H^p(X, \mathcal{F}') = 0$ for all $p \geq 1$, and
\item $\{H^n(X, -)\}_{n \geq 0}$
is a universal $\delta$-functor from $\QCoh(\mathcal{O}_X)$ to
$\textit{Ab}$.
\end{enumerate}
\end{lemma}

\begin{proof}
Let $X = \bigcup U_i$ be a finite affine open covering.
Set $U = \coprod U_i$ and denote $j : U \to X$
the morphism inducing the given open immersions $U_i \to X$.
Since $U$ is an affine scheme and $X$ has affine diagonal,
the morphism $j$ is affine, see
Morphisms, Lemma \ref{morphisms-lemma-affine-permanence}.
For every $\mathcal{O}_X$-module $\mathcal{F}$ there is
a canonical map $\mathcal{F} \to j_*j^*\mathcal{F}$.
This map is injective as can be seen by checking on stalks:
if $x \in U_i$, then we have a factorization
$$
\mathcal{F}_x \to (j_*j^*\mathcal{F})_x
\to (j^*\mathcal{F})_{x'} = \mathcal{F}_x
$$
where $x' \in U$ is the point $x$ viewed as a point of $U_i \subset U$.
Now if $\mathcal{F}$ is quasi-coherent, then $j^*\mathcal{F}$
is quasi-coherent on the affine scheme $U$ hence has vanishing
higher cohomology by
Lemma \ref{lemma-quasi-coherent-affine-cohomology-zero}.
Then $H^p(X, j_*j^*\mathcal{F}) = 0$ for
$p > 0$ by Lemma \ref{lemma-relative-affine-cohomology}
as $j$ is affine. This proves (1).
Finally, we see that the map
$H^p(X, \mathcal{F}) \to H^p(X, j_*j^*\mathcal{F})$
is zero and part (2) follows from
Homology, Lemma \ref{homology-lemma-efface-implies-universal}.
\end{proof}

\begin{lemma}
\label{lemma-vanishing-nr-affines-quasi-separated}
Let $X$ be a quasi-compact quasi-separated scheme.
Let $X = U_1 \cup \ldots \cup U_n$ be an open covering
with each $U_i$ quasi-compact and separated (for example affine).
Set
$$
d = \max\nolimits_{I \subset \{1, \ldots, n\}}
\left(|I| + t(\bigcap\nolimits_{i \in I} U_i) - 1\right)
$$
where $t(U)$ is the minimal number of affines needed to cover
the scheme $U$. Then $H^p(X, \mathcal{F}) = 0$ for all $p \geq d$ and all
quasi-coherent sheaves $\mathcal{F}$.
\end{lemma}

\begin{proof}
Note that since $X$ is quasi-separated and $U_i$ quasi-compact the numbers
$t(\bigcap_{i \in I} U_i)$ are finite. Proof using induction on $n$.
If $n = 1$ then the result follows from Lemma \ref{lemma-vanishing-nr-affines}.
If $n > 1$, write $X = U \cup V$ with $U = U_1 \cup \ldots \cup U_{n - 1}$
and $V = U_n$. We apply the Mayer-Vietoris long exact sequence
$$
0 \to
H^0(X, \mathcal{F}) \to
H^0(U, \mathcal{F}) \oplus H^0(V, \mathcal{F}) \to
H^0(U \cap V, \mathcal{F}) \to
H^1(X, \mathcal{F}) \to \ldots
$$
see Cohomology, Lemma \ref{cohomology-lemma-mayer-vietoris}.
To finish the proof for $q \geq d$ we will show that
$H^q(V, \mathcal{F})$, $H^q(U, \mathcal{F})$, and
$H^{q - 1}(U \cap V, \mathcal{F})$ vanish.
By the case $n = 1$ we have $H^q(V, \mathcal{F}) = 0$ for
$q \geq t(V) = t(U_n)$. Since $t(V) \leq d$ this proves what we want.
By induction hypothesis we have $H^q(U, \mathcal{F}) = 0$ for
$$
q \geq \max\nolimits_{I \subset \{1, \ldots, n - 1\}}
\left(|I| + t(\bigcap\nolimits_{i \in I} U_i) - 1\right)
$$
Since the integer on the right is less than or equal to $d$, this
proves what we want. Finally we may use our induction hypothesis for the open
$U \cap V = (U_1 \cap U_n) \cup \ldots \cup (U_{n - 1} \cap U_n)$ to get
the vanishing of $H^q(U \cap V, \mathcal{F}) = 0$ for
$$
q \geq \max\nolimits_{I \subset \{1, \ldots, n - 1\}}
\left(|I| + t(U_n \cap \bigcap\nolimits_{i \in I} U_i) - 1\right)
$$
Since the integer on the right is strictly less than $d$
the lemma follows.
\end{proof}

\begin{lemma}
\label{lemma-quasi-coherence-higher-direct-images}
Let $f : X \to S$ be a morphism of schemes.
Assume that $f$ is quasi-separated and quasi-compact.
\begin{enumerate}
\item For any quasi-coherent $\mathcal{O}_X$-module $\mathcal{F}$ the
higher direct images $R^pf_*\mathcal{F}$ are quasi-coherent on $S$.
\item If $S$ is quasi-compact, there exists an integer $n = n(X, S, f)$
such that $R^pf_*\mathcal{F} = 0$ for all $p \geq n$ and any
quasi-coherent sheaf $\mathcal{F}$ on $X$.
\item In fact, if $S$ is quasi-compact we can find $n = n(X, S, f)$
such that for every
morphism of schemes $S' \to S$ we have $R^p(f')_*\mathcal{F}' = 0$
for $p \geq n$ and any quasi-coherent sheaf $\mathcal{F}'$
on $X'$. Here $f' : X' = S' \times_S X \to S'$ is the base change of $f$.
\end{enumerate}
\end{lemma}

\begin{proof}
We first prove (1). Note that under the hypotheses of the lemma the sheaf
$R^0f_*\mathcal{F} = f_*\mathcal{F}$ is quasi-coherent by
Schemes, Lemma \ref{schemes-lemma-push-forward-quasi-coherent}.
Using
Cohomology, Lemma \ref{cohomology-lemma-localize-higher-direct-images}
we see that forming higher direct images commutes with restriction
to open subschemes. Since being quasi-coherent is local on $S$ we
reduce to the case discussed in the next paragraph.

\medskip\noindent
Proof of (1) in case $S$ is affine. We will use the induction principle.
Since $f$ quasi-compact and quasi-separated we see that $X$
is quasi-compact and quasi-separated. For $U \subset X$
quasi-compact open and $a = f|_U$ we let $P(U)$ be the property that
$R^pa_*\mathcal{F}$ is quasi-coherent on $S$ for all quasi-coherent modules
$\mathcal{F}$ on $U$ and all $p \geq 0$. Since $P(X)$ is (1), it suffices
the prove conditions (1) and (2) of Lemma \ref{lemma-induction-principle} hold.
If $U$ is affine, then $P(U)$ holds because $R^pa_*\mathcal{F} = 0$ for
$p \geq 1$ (by  Lemma \ref{lemma-relative-affine-vanishing} and
Morphisms, Lemma \ref{morphisms-lemma-morphism-affines-affine})
and we've already observed the result holds for $p = 0$ in the first
paragraph. Next, let $U \subset X$ be a quasi-compact open, $V \subset X$
an affine open, and assume $P(U)$, $P(V)$, $P(U \cap V)$ hold.
Let $a = f|_U$, $b = f|_V$, $c = f|_{U \cap V}$, and $g = f|_{U \cup V}$.
Then for any quasi-coherent $\mathcal{O}_{U \cup V}$-module $\mathcal{F}$
we have the relative Mayer-Vietoris sequence
$$
0 \to
g_*\mathcal{F} \to
a_*(\mathcal{F}|_U) \oplus b_*(\mathcal{F}|_V) \to
c_*(\mathcal{F}|_{U \cap V}) \to
R^1g_*\mathcal{F} \to \ldots
$$
see Cohomology, Lemma \ref{cohomology-lemma-relative-mayer-vietoris}.
By $P(U)$, $P(V)$, $P(U \cap V)$ we see that
$R^pa_*(\mathcal{F}|_U)$, $R^pb_*(\mathcal{F}|_V)$ and
$R^pc_*(\mathcal{F}|_{U \cap V})$ are all quasi-coherent.
Using the results on quasi-coherent sheaves in
Schemes, Section \ref{schemes-section-quasi-coherent}
this implies that each of the sheaves
$R^pg_*\mathcal{F}$ is quasi-coherent since it sits in the middle of a short
exact sequence with a cokernel of a map between quasi-coherent sheaves
on the left and a kernel of a map between quasi-coherent sheaves on the right.
Whence $P(U \cup V)$ and the proof of (1) is complete.

\medskip\noindent
Next, we prove (3) and a fortiori (2). Choose a finite affine open
covering $S = \bigcup_{j = 1, \ldots m} S_j$. For each $j$ choose
a finite affine open covering
$f^{-1}(S_j) = \bigcup_{i = 1, \ldots t_j} U_{ji} $.
Let
$$
d_j = \max\nolimits_{I \subset \{1, \ldots, t_j\}}
\left(|I| + t(\bigcap\nolimits_{i \in I} U_{ji})\right)
$$
be the integer found in
Lemma \ref{lemma-vanishing-nr-affines-quasi-separated}.
We claim that $n(X, S, f) = \max d_j$ works.

\medskip\noindent
Namely, let $S' \to S$ be a morphism of schemes and let
$\mathcal{F}'$ be a quasi-coherent sheaf on $X' = S' \times_S X$.
We want to show that $R^pf'_*\mathcal{F}' = 0$ for $p \geq n(X, S, f)$.
Since this question is local on $S'$ we may assume that $S'$ is affine
and maps into $S_j$ for some $j$. Then $X' = S' \times_{S_j} f^{-1}(S_j)$
is covered by the open affines $S' \times_{S_j} U_{ji}$, $i = 1, \ldots t_j$
and the intersections
$$
\bigcap\nolimits_{i \in I} S' \times_{S_j} U_{ji} =
S' \times_{S_j} \bigcap\nolimits_{i \in I} U_{ji}
$$
are covered by the same number of affines as before the base change.
Applying
Lemma \ref{lemma-vanishing-nr-affines-quasi-separated}
we get $H^p(X', \mathcal{F}') = 0$. By the first part of the proof
we already know that each $R^qf'_*\mathcal{F}'$ is quasi-coherent
hence has vanishing higher cohomology groups on our affine scheme $S'$,
thus we see that $H^0(S', R^pf'_*\mathcal{F}') = H^p(X', \mathcal{F}') = 0$
by Cohomology, Lemma \ref{cohomology-lemma-apply-Leray}.
Since $R^pf'_*\mathcal{F}'$ is quasi-coherent
we conclude that $R^pf'_*\mathcal{F}' = 0$.
\end{proof}

\begin{lemma}
\label{lemma-quasi-coherence-higher-direct-images-application}
Let $f : X \to S$ be a morphism of schemes.
Assume that $f$ is quasi-separated and quasi-compact.
Assume $S$ is affine.
For any quasi-coherent $\mathcal{O}_X$-module $\mathcal{F}$
we have
$$
H^q(X, \mathcal{F}) = H^0(S, R^qf_*\mathcal{F})
$$
for all $q \in \mathbf{Z}$.
\end{lemma}

\begin{proof}
Consider the Leray spectral sequence $E_2^{p, q} = H^p(S, R^qf_*\mathcal{F})$
converging to $H^{p + q}(X, \mathcal{F})$, see
Cohomology, Lemma \ref{cohomology-lemma-Leray}.
By Lemma \ref{lemma-quasi-coherence-higher-direct-images}
we see that the sheaves $R^qf_*\mathcal{F}$ are quasi-coherent.
By Lemma \ref{lemma-quasi-coherent-affine-cohomology-zero}
we see that $E_2^{p, q} = 0$ when $p > 0$.
Hence the spectral sequence degenerates at $E_2$ and we win.
See also
Cohomology, Lemma \ref{cohomology-lemma-apply-Leray} (2)
for the general principle.
\end{proof}








